\subsection{A divisorial calculation for the last Cartan layer}\label{sec:divisorial-cartan}
This subsection computes the ramification and inseparable residue extension of a concrete quotient of the universal deformation surface. Its complete discrete valuation fields have imperfect residue field, so the calculation is a supporting result rather than an example in the standing perfect-residue-field class.

Let $k=\overline{\F}_p$, and let $R_m$ be the complete local ring of the universal deformation of a one-dimensional height-two formal group over $k$, with full Drinfeld level-$p^m$ structure. Write $G$ for the universal formal group over $R_0$. For $m\ge1$, write $x_m,y_m$ for the coordinates of a chosen level basis. These are regular parameters of the two-dimensional ring $R_m$, and, for $m\ge2$, the transition is
\[
 x_{m-1}=[p](x_m),\qquad y_{m-1}=[p](y_m).
\]
The level action of $\GL_2(\Z/p^m)$ fixes the deformation base $R_0$ and $W(k)$. These are Drinfeld's regularity and level-parameter results, in the form of \citet[Theorem~2.1.2 and its proof, pp.~9--11; \S2.2, pp.~12--14]{Strauch2006}; see also \citet[Theorem~2.1.2, p.~437; Proposition~2.2.3(i), p.~439]{CarterStrauch2021}. The height assumption gives a nonzero degree-$p^2$ leading term of the special $p$-series. Scaling the coordinate over $k$ normalizes its coefficient to one. No exact polynomial expression for a mixed-characteristic lift is required.

For comparison, \citet[Section~2.2.8 and Corollary~2.2.9, p.~443]{CarterStrauch2021} obtain ramification index $(p-1)p^{n-1}$, when $\mathcal O=\Z_p$, for a tower formed by localizing at the height-one prime generated by one level coordinate. The valuation used here assigns positive weight to both level parameters, so its scalar fixed fields require a separate calculation.

\begin{proposition}\label{prop:divisorial-last-cartan}
Let $p$ be odd and $n\ge2$. Put
\[
 A=p^2+p+1,\qquad B=p^2,\qquad N=p^{2n-2}.
\]
Let $w$ be the divisorial monomial valuation on $\operatorname{Frac}(R_n)$ with $w(x_n)=A$ and $w(y_n)=B$, normalized to value group $\Z$, and let $F$ be its completion. Then its residue field is $k(t)$, where
\[
 t=\operatorname{res}(x_n^B/y_n^A),
\]
and
\begin{equation}\label{eq:divisorial-level-values}
 w(x_j)=A p^{2n-2j},\qquad w(y_j)=B p^{2n-2j}
 \quad(1\le j\le n),\qquad
 w(p)=(p^4-1)N.
\end{equation}
Fix a nonsquare $\epsilon\in\Z_p^\times$ and put
\[
 J=\begin{pmatrix}0&\epsilon\\1&0\end{pmatrix},\qquad
 U=I+p^{n-1}I,\qquad V=I+p^{n-1}J,\qquad H=\langle U,V\rangle .
\]
The level actions preserve $w$ and extend faithfully to $F$. The elements $U,V$ have order $p$, and $H$ has order $p^2$. The exact relative invariants are
\[
\begin{array}{c|cc}
 \text{extension}&\text{ramification index}&\text{residue degree}\\ \hline
 F/F^{\langle U\rangle}&1&p\\
 F/F^{\langle V\rangle}&p&1\\
 F/F^H&p&p .
\end{array}
\]
Each residue extension is purely inseparable, with the trivial extension included in this terminology. The full principal Cartan subgroup
\[
 C=(1+p\Z_p[J])/(1+p^n\Z_p[J])
\]
also preserves $w$, has order $p^{2n-2}$, and satisfies
\begin{equation}\label{eq:divisorial-full-cartan-bound}
 f(F/F^C)\ge p,\qquad e(F/F^C)\le p^{2n-3}.
\end{equation}
No equality for the full subgroup $C$ is asserted.
\end{proposition}

\begin{proof}
\emph{The valuation and the torsion values.}
The monomial valuation exists over the mixed-characteristic regular surface without choosing a coefficient field $k$ inside $R_n$. Blow up the closed point repeatedly. When the positive weights are $a>b$, the chart with parameters $x/y,y$ has weights $a-b,b$; use the other chart when $b>a$. The Euclidean algorithm reaches weights $1,1$, and the generic point of the next exceptional divisor is the required discrete valuation ring. Reversing these unimodular monomial changes gives residue field $k(t)$ and value group $\Z$. If $X,Y$ denote the initial forms of $x=x_n,y=y_n$, the associated graded fraction field is
\begin{equation}\label{eq:divisorial-graded-field}
 \operatorname{Frac}(\operatorname{gr}_w F)=k(X,Y).
\end{equation}
Indeed, choose integers $a,b$ with $aA+bB=1$. The monomials $x^ay^b$ and $x^B/y^A$ give a uniformizer and a transcendental residue coordinate, and their exponent matrix has determinant $-1$.

Let $m$ be the minimum value of a nonzero $p$-torsion point. The full-level product gives
\[
 [p](T)=q(T)\prod_{P\in G[p]}(T-T(P)),
\]
where $q$ is a unit whose constant residue is one. If $0<w(z)<m$, every factor has initial form $z$, including the factor for the zero point. Thus
\begin{equation}\label{eq:divisorial-small-root}
 w([p]z)=p^2w(z),\qquad
 \operatorname{in}_w([p]z)=\operatorname{in}_w(z)^{p^2}.
\end{equation}
One of $x_1,y_1$ attains $m$: a formal $\F_p$-linear combination cannot have value smaller than both basis points.

Multiplication by $p$ strictly raises the value of every positive-valued point with nonzero image, because its linear term is $pT$ and every other term has degree at least two. All ancestors of a minimum point therefore have value below $m$. Formula~\eqref{eq:divisorial-small-root} shows that either $m=NA$ or $m=NB$. The first possibility is impossible: starting from $y$ of value $B<A$, iteration would give $w(y_1)=NB<NA=m$. Hence $m=NB$ and the entire $y$-chain has the displayed values. In the $x$-chain the last input before reaching $x_1$ has value $AN/p^2<BN$. The same formula proves all its values and, for both chains,
\[
 \operatorname{in}_w(x_j)=X^{p^{2n-2j}},
 \qquad \operatorname{in}_w(y_j)=Y^{p^{2n-2j}}.
\]
At the first level, $p-1$ nonzero points have value $AN$, and the other $p^2-p$ have value $BN$. The linear coefficient of the product is $p$ up to a unit. Therefore
\[
 w(p)=\bigl((p-1)A+(p^2-p)B\bigr)N=(p^4-1)N,
\]
which proves~\eqref{eq:divisorial-level-values}.

\emph{The action and its leading derivations.}
An element of $C$ has the form $aI+bJ$, with $a\equiv1\pmod p$ and $b\equiv0\pmod p$. Its off-diagonal contribution to $x$ has value at least $p^2B>A$, and the correction to its diagonal contribution has value at least $p^2A>A$. The corresponding corrections to $y$ also have greater value than $y$. Thus the initial forms of $x,y$ are fixed. The action fixes $W(k)$, preserves $w$, and extends to $F$. It remains faithful because the original fraction field embeds in its completion. The two last-layer matrices have order $p$ and commute. Their coefficients $I,J$ in the last congruence layer are linearly independent over $\F_p$, giving $|H|=p^2$.

Use the basis-column convention $\phi\mapsto\phi\circ g$. Formal addition and the torsion initial forms give
\[
\begin{aligned}
 \operatorname{in}_w(Ux-x)&=X^N,&
 \operatorname{in}_w(Uy-y)&=Y^N,\\
 \operatorname{in}_w(Vx-x)&=Y^N,&
 \operatorname{in}_w(Vy-y)&=\overline\epsilon X^N .
\end{aligned}
\]
A transposed level convention moves the nonzero scalar $\overline\epsilon$ between the last two formulas and has no effect on the argument.
The smaller displacement excess for $U$ is $B(N-1)$, attained on $y$; for $V$ it is $BN-A$, attained on $x$. Both are positive. The first nonzero graded differences therefore define derivations of~\eqref{eq:divisorial-graded-field}:
\begin{equation}\label{eq:divisorial-derivations}
 D_U=Y^N\frac{\partial}{\partial Y},\qquad
 D_V=Y^N\frac{\partial}{\partial X}.
\end{equation}

Here the derivation assertion applies to the complete field. To justify it, first work on $R_n$: lift a homogeneous initial form by a polynomial in $x,y$ with Teichm\"uller coefficients from $W(k)$, subtract, and repeat. The positive monomial weights give convergent approximations; this filtration is comparable with the maximal-ideal filtration. The action fixes the coefficient lifts. If $d>0$ is the smaller displacement excess, the identity
\[
 g(fh)-fh=(g(f)-f)h+f(g(h)-h)+(g(f)-f)(g(h)-h)
\]
shows that the degree-$d$ term is a derivation, since the last summand has larger degree. The bound and its initial term pass to the convergent approximations, to fractions by the quotient rule, and then to the completion. In particular a higher-weight tail cannot cancel a nonzero value of the derivation in~\eqref{eq:divisorial-derivations}.

\emph{Fixed graded fields.}
For a finite group $\Gamma$ of these automorphisms, $E=F^\Gamma$ is complete and $F/E$ is Galois of degree $|\Gamma|$. A finite extension of complete discrete valuation fields satisfies $[F:E]=ef$, with the total residue degree, even if that residue extension is inseparable. One way to see this is to use the finite integral closure $\Ocal_F$ over $\Ocal_E$: it is free of rank $[F:E]$, and reduction modulo a uniformizer of $E$ has length $ef$. Finiteness follows by placing integral elements in the trace-dual of a lattice generated by an integral field basis; separability makes that trace pairing nondegenerate. The same valuation classes and a residue-field basis give
\begin{equation}\label{eq:divisorial-graded-degree}
 [k(X,Y):\operatorname{Frac}(\operatorname{gr}_w E)]=ef=|\Gamma|.
\end{equation}
Independence is checked first in distinct valuation classes modulo $e$, then in the residue field; successive leading-term subtraction proves spanning.

An element fixed by $U$ has initial form killed by $D_U$, so
\[
 \operatorname{Frac}(\operatorname{gr}_w F^{\langle U\rangle})
 \subset k(X,Y^p).
\]
Both fields on the right and left have degree $p$ in $k(X,Y)$, by~\eqref{eq:divisorial-graded-degree}; hence equality holds. The weight group of $k(X,Y^p)$ is generated by $A,pB$ and is $\Z$. Its degree-zero field is $k(t^p)$. This gives $e=1,f=p$ for the scalar subgroup.

For $V$ the same argument gives
\[
 \operatorname{Frac}(\operatorname{gr}_w F^{\langle V\rangle})
 =k(X^p,Y).
\]
Its weight group is $p\Z$, since $\gcd(pA,B)=p$, and its degree-zero field is $k(t)$, since $(X^p)^p/Y^A=X^B/Y^A$. This gives $e=p,f=1$.

For $H$, both derivations kill every invariant initial form. Their common kernel is $k(X^p,Y^p)$, because $k$ is perfect. Its degree is $p^2$, so~\eqref{eq:divisorial-graded-degree} again gives equality with the fixed graded fraction field. Its weight group is $p\Z$ and its degree-zero field is $k(t^p)$. Thus $e=f=p$ for $H$, proving the table and its purely inseparable residue assertions.

\emph{The full principal subgroup.}
The principal-unit filtration of the unramified quadratic ring $\Z_p[J]$ gives $|C|=p^{2n-2}$. Since $F^C\subset F^H$, residue degrees in the tower multiply, giving $f(F/F^C)\ge p$. The degree formula gives~\eqref{eq:divisorial-full-cartan-bound}. This does not compute that full residue degree.
\end{proof}

The last-layer bound alone gives a weaker numerical observation at $p=5,n=4$. If $F^C/K$ is a tame extension of ramification index $12$ with unchanged residue field, then
\[
 e_K=\frac{624\cdot5^6}{12\,e(F/F^C)}
 \ge\frac{624\cdot5^6}{12\cdot5^5}=260.
\]
Equality in this estimate would require $e(F/F^C)=5^5$, or equivalently $f(F/F^C)=5$. The next proposition excludes that equality for this valuation by computing the full scalar subgroup. The tame extension remains an additional hypothesis, and no perfect-residue specialization follows from the divisorial calculation.
